% 例 3 三角形内角和：过 A 作 BC 平行线，∠A + ∠B + ∠C = 平角
\documentclass[tikz,border=4pt]{standalone}
\usepackage{ctex}
\usepackage{amsmath}
\usetikzlibrary{calc, arrows.meta, decorations.markings}
\begin{document}
\begin{tikzpicture}[scale=1.0, thick,
  arrowone/.style={postaction={decorate, decoration={markings,
    mark=at position 0.5 with {\arrow{Stealth}}}}}]
  \coordinate (A) at (2,3);
  \coordinate (B) at (0,0);
  \coordinate (C) at (5,0);
  \draw[arrowone] (B) -- (C);
  \draw (A) -- (B);
  \draw (A) -- (C);
  % 过 A 作平行线
  \draw[red, dashed, arrowone] (-1,3) -- (6,3);
  % 角
  \pgfmathsetmacro{\angAB}{atan2(0-3, 0-2)} % from A to B
  \pgfmathsetmacro{\angAC}{atan2(0-3, 5-2)} % from A to C
  % ∠B at B
  \draw (B) ++(0.6,0) arc (0:atan2(3,2):0.6);
  \node[font=\footnotesize] at ($(B)+(0.55,0.3)$) {$\angle B$};
  % ∠C at C
  \draw (C) ++(-0.6,0) arc (180:atan2(3,-3):0.6);
  \node[font=\footnotesize] at ($(C)+(-0.55,0.3)$) {$\angle C$};
  % ∠A at A
  \draw (A) ++({0.6*cos(\angAB)},{0.6*sin(\angAB)}) arc (\angAB:\angAC:0.6);
  \node[font=\footnotesize] at ($(A)+(0,-0.4)$) {$\angle A$};
  % 在 A 处沿平行线两侧的内错角：左侧 = ∠B，右侧 = ∠C
  \draw[blue] (A) ++(-0.6,0) arc (180:\angAB:0.6);
  \node[blue, font=\footnotesize] at ($(A)+(-0.5,0.25)$) {$\angle B$};
  \draw[blue] (A) ++(0.6,0) arc (0:\angAC:0.6);
  \node[blue, font=\footnotesize] at ($(A)+(0.5,0.25)$) {$\angle C$};
  \foreach \pt/\pos in {A/above, B/below left, C/below right}
    \filldraw (\pt) circle (1.3pt) node[\pos] {$\pt$};
  \node[below, font=\footnotesize] at (2.5,-0.6) {$\angle A + \angle B + \angle C = 180°$（平角）};
\end{tikzpicture}
\end{document}
